\subsection{Rules sufficient for elementary models}

The computational goal is deliberately limited. We need enough rules to differentiate low-degree polynomials, exponentials, sine and cosine, and simple compositions.

\begin{theorembox}[Basic rules]
For differentiable functions,
\[
(c)'=0,
\qquad
(x^n)'=nx^{n-1},
\qquad
(cf)'=cf',
\qquad
(f+g)'=f'+g'.
\]
We also use
\[
(e^x)'=e^x,
\qquad
(\sin x)'=\cos x,
\qquad
(\cos x)'=-\sin x.
\]
\end{theorembox}

\begin{examplebox}[A polynomial]
For
\[
f(x)=x^3-3x,
\]
we obtain
\[
f'(x)=3x^2-3.
\]
\end{examplebox}

\begin{examplebox}[A simple oscillation]
If
\[
y(t)=2\sin t,
\]
then
\[
y'(t)=2\cos t.
\]
The derivative describes how rapidly the oscillating quantity changes.
\end{examplebox}

\begin{theorembox}[Product, quotient, and chain rules]
\[
(fg)'=f'g+fg',
\qquad
\left(\frac{f}{g}\right)'=\frac{f'g-fg'}{g^2},
\]
and
\[
(f\circ g)'(x)=f'(g(x))g'(x).
\]
\end{theorembox}

\begin{examplebox}[An elementary chain rule]
For
\[
f(x)=(x^2+1)^3,
\]
\[
f'(x)=3(x^2+1)^2\cdot2x.
\]
The example is chosen to display the nested structure, not to create difficult algebra.
\end{examplebox}
